\section{Introduction}
\label{sec:intro}

Modern observational facilities have ushered in an era of three-dimensional spectroscopy at an unprecedented scale. Integral field unit (IFU) spectrographs such as MUSE on the Very Large Telescope \citep{Bacon2010} and millimeter/submillimeter interferometers such as ALMA now routinely deliver spectral data cubes in which every spatial pixel carries a complete spectrum across hundreds to thousands of velocity channels. These cubes are not merely images with attached spectra: each physically distinct structure --- a rotating gaseous disk, an AGN-driven outflow, a tidal bridge connecting merging companions --- traces a coherent locus through the full three-dimensional $(x, y, v)$ space of the cube, encoding kinematic, morphological, and energetic information that is irretrievably compressed when the cube is collapsed into a two-dimensional projection.

Despite this richness, source extraction in radio and millimeter astronomy has remained predominantly a two-dimensional enterprise. The canonical workflow collapses spectral cubes into moment-0 (velocity-integrated intensity) maps and then applies tools such as \textsc{SExtractor} \citep{Bertin1996} or \textsc{PyBDSF} \citep{Mohan2015} --- software designed for continuum imaging that treats the third axis of the cube as a nuisance to be summed away rather than a dimension to be exploited. The consequences of this mismatch are concrete. Spectrally narrow sources confined to only a few velocity channels are heavily diluted below detection thresholds in broadband moment maps, while bright nearby sources contaminate their footprints. Morphologically complex systems --- galaxy pairs, objects undergoing major mergers, host galaxies of actively accreting supermassive black holes --- produce emission that shifts, stretches, and bifurcates across velocity channels; a two-dimensional finder sees only their aggregate projection and cannot recover kinematically distinct components. Most critically, the spatial centroid of a rotating structure can shift by several synthesized beams across the velocity axis \citep[e.g.,][]{Venemans2020, Decarli2021}, meaning that even a clean detection on the moment-0 map misrepresents the source position at any individual velocity. For science that requires accurately resolved kinematic maps --- rotation curves, outflow opening angles, star formation rate surface densities --- these positional and morphological errors propagate directly into derived physical quantities.

Genuinely three-dimensional source finders do exist, and their development has been driven primarily by neutral hydrogen (HI) 21\,cm spectral-line surveys. \textsc{Duchamp} \citep{Whiting2012} detects sources by thresholding in three dimensions and linking connected voxels, offering a conceptual advance over moment-map approaches in that spectral coherence is a detection criterion. \textsc{SoFiA} \citep{Serra2015} and its parallelised successor \textsc{SoFiA-2} \citep{Westmeier2021} extended this framework with multi-scale spatial and spectral smoothing kernels, reliability estimation from negative detections, and highly optimised pipelines capable of processing the large data volumes anticipated from ASKAP/WALLABY. These tools have proven enormously productive for low-redshift HI surveys, achieving near-complete source recovery above integrated signal-to-noise ratios of $\sim5$--6 \citep{Westmeier2021}. A complementary approach is offered by hierarchical dendrogram decomposition \citep[as implemented in \textsc{astrodendro};][]{Rosolowsky2008}, which segments the emission field into a nested tree of structures --- leaves, branches, and trunks --- defined by iso-intensity contours. \textsc{astrodendro} has found broad application in the study of molecular cloud hierarchies and high-density gas tracers, where the physical goal is to characterise the full hierarchy of sub-structures rather than to detect point-like sources above a noise floor.

However, none of these tools was designed for the specific challenge this paper addresses: the detection and tracking of \emph{kinematically offset}, \emph{morphologically non-stationary} sources in high-resolution interferometric data of distant, high-surface-brightness systems. Several features of this regime combine to make existing tools insufficient. First, in interferometric cubes the noise is spatially correlated by the synthesised beam, and its spectral covariance structure depends on the frequency-dependent primary beam and baseline distribution \citep{Loomis2018}; the isotropic Gaussian noise model assumed by threshold-based finders is a poor approximation. Second, the objects of primary scientific interest at high redshift --- massive star-forming galaxies, quasar host galaxies, submillimetre galaxies --- are neither spatially unresolved point sources nor smoothly extended low-surface-brightness discs: they exhibit steep velocity gradients, irregular morphologies, and companions at small projected separations that can overlap both spatially and spectrally. Third, and most fundamentally, the spatial distribution of emission from a rotating or interacting system is not stationary in the spectral dimension. A source that appears as an elongated arc on the receding side of a disk at one velocity channel will appear as a different arc --- shifted, rotated, or split --- at the next channel. Finders that link detections via fixed apertures or connected-component labelling, without any model for \emph{how} emission is expected to move between channels, cannot distinguish this kinematic drift from noise fluctuations or source confusion.

This problem is particularly acute for observations of high-redshift luminous systems in fine-structure and molecular emission lines. At $z \sim 4$--7, the [C\,\textsc{ii}] $158\,\mu$m line is one of the brightest ISM diagnostics accessible to ALMA \citep{Carilli2013}, and it has become the workhorse tracer for resolved kinematics of early-universe star-forming galaxies and quasar hosts \citep{Venemans2020, Decarli2021}. High-resolution [C\,\textsc{ii}] cubes of hyperluminous infrared galaxies (HyLIRGs) and quasar host--companion pairs reveal velocity fields of remarkable complexity: disks with steep rotation curves whose centroid drifts by $\gtrsim 1\arcsec$ across the line profile \citep{Venemans2020}, companions interacting tidally whose [C\,\textsc{ii}] emission stretches coherently over $>25\,\mathrm{kpc}$ and $>1500\,\mathrm{km\,s}^{-1}$ in projected phase space \citep{Decarli2021}, and large-scale molecular outflows whose position angle and radial extent change channel by channel \citep{Maiolino2012}. In such data, the question is not merely ``is there a source?'' but ``what is the three-dimensional trajectory of this source through the cube, and does that trajectory correspond to a coherent physical object?''

We introduce NEMO (\textbf{N}on-stationary \textbf{E}xtraction via \textbf{M}ultiscale \textbf{O}ptical-flow), a source finder built around a fundamentally different detection philosophy. NEMO proceeds in three coupled stages. First, it detects emission candidates in individual velocity channels using a noise model calibrated to interferometric data, exploiting matched filtering \citep{Loomis2018} to improve sensitivity in the low-SNR, beam-correlated regime. Second, it estimates the local velocity field between consecutive channel pairs using dense optical flow \citep{Horn1981, Lucas1981} --- a technique from computer vision that computes, for every spatial pixel, the apparent displacement field that best explains the change in intensity between two frames. Applied to a spectral cube, optical flow produces a per-pixel kinematic map that captures the continuous drift of emitting structures across channels without requiring any prior knowledge of the source geometry. Third, NEMO uses the estimated flow field as a propagation model to link per-channel detections into coherent three-dimensional tracks, and applies a reliability filter analogous to that of \citet{Serra2015} to reject spurious linkages. The output is a catalogue of sources defined not by a static aperture but by a dynamic trajectory through $(x, y, v)$ space, together with per-source moment maps, integrated spectra, and kinematic summary statistics.

This paper is organised as follows. Section~\ref{sec:data} describes the ALMA [C\,\textsc{ii}] and dust-continuum data used to validate NEMO. Section~\ref{sec:method} details the algorithm --- the noise model, the matched filter, the optical-flow estimation, and the track-linking procedure. Section~\ref{sec:results} presents detection completeness and purity tests using injected mock sources, and compares NEMO's output to \textsc{SoFiA-2} and \textsc{astrodendro} on the same cubes. Section~\ref{sec:science} demonstrates NEMO's scientific utility on the target sample. Section~\ref{sec:conclusions} summarises our findings.

